\section{Related Work}

This paper sits at the intersection of retrieval-augmented generation, verifier-and-repair methods, agentic workflow design, and safety-sensitive advisory systems. We organize the related work around four questions: how generated claims are grounded in evidence, how critique or repair reduces unsupported generation, how tool-using workflows expose intermediate decisions, and what changes when the output is advice in a risk-sensitive setting. The gap is not the absence of RAG, agents, or self-critique. The gap is that refusal, bounded repair, and safety-boundary behavior are rarely treated as explicit measurable workflow states in a focused advisory benchmark.

\subsection{RAG Grounding and Attribution}

Retrieval-augmented generation connects parametric generation with non-parametric retrieved evidence for knowledge-intensive tasks \cite{lewis2020rag}. A major follow-up concern is whether generated answers are actually supported by the sources they cite. Work on generative search verifiability evaluates whether claims can be checked against cited sources \cite{liu2023verifiability}; citation-generation work studies how to produce long-form answers with source attributions \cite{gao2023alce}; RAGTruth annotates hallucinations in RAG outputs \cite{niu2023ragtruth}; and RAGCHECKER proposes fine-grained diagnostics for retriever and generator behavior \cite{ragchecker2024}. This line of work establishes that retrieval and citation are not enough: a generated answer can still contain unsupported claims, incomplete support, or misleading attribution.

Our setting uses this insight but changes the diagnostic target. In safety-sensitive endurance-training advice, the relevant question is not only whether a final answer is factually supported, but also whether the workflow can expose when evidence is sufficient, partial, missing, conflicting, or unsafe to act on. For this reason, the benchmark distinguishes \texttt{answered}, \texttt{partial\_answer}, and \texttt{refused}, and treats designed-unanswerable refusal as an intended diagnostic outcome rather than as a generic non-answer.

\subsection{Self-Critique, Verification, and Repair}

Recent retrieval and verification methods make critique more explicit. Self-RAG learns to retrieve, generate, and critique through self-reflection \cite{asai2023selfrag}, while Corrective RAG evaluates retrieved documents and adjusts generation when retrieval quality is low \cite{yan2024crag}. Chain-of-Verification reduces hallucination by drafting, verifying, and revising model outputs \cite{dhuliawala2023cove}. Reflexion and Self-Refine study iterative improvement through verbal feedback or self-feedback \cite{shinn2023reflexion,madaan2023selfrefine}. Together, these methods show that critique and revision can improve factuality or task performance beyond one-shot generation.

The limitation for our purpose is that critique and repair are often evaluated as mechanisms for improving the final answer, while the intermediate failure modes remain less visible. Our workflow therefore separates Evidence Gate, Risk Gate, Independent Grounding Auditor, and Bounded Citation/Grounding Repair as logged control points. This makes \texttt{citation repair}, \texttt{false refusal}, unanswerable-item refusal, and request-level blocking countable workflow states rather than post-hoc explanations of a single final response. This is a diagnostic use of repair, not a claim that repair alone makes the advice safe.

\subsection{Agentic Workflow and Tool Use}

Tool-using and modular-agent systems illustrate how LLM applications can be decomposed into explicit reasoning, acting, memory, and tool-use stages. ReAct interleaves reasoning and acting \cite{yao2022react}; MRKL routes LLMs to external knowledge and symbolic modules \cite{karpas2022mrkl}; Toolformer studies learned tool use \cite{schick2023toolformer}; ReWOO decouples reasoning from observations \cite{xu2023rewoo}; LLM Compiler explores parallel function calling \cite{kim2023llmcompiler}; and MemGPT frames memory management as an operating-system-like concern for LLM agents \cite{packer2023memgpt}. Multi-agent frameworks further show that LLM systems can assign roles and coordinate specialized agents for complex tasks \cite{wu2023autogen,hong2023metagpt,du2023debate}.

These systems motivate the engineering shape of our method, but they also define the boundary of the contribution. We do not claim that modularity, tool use, or multi-agent orchestration is new. Instead, we use the workflow perspective to expose a narrower set of control points for advisory RAG: evidence validation, risk gating, evidence-constrained generation, independent grounding audit, bounded repair, and refusal. The contribution is not a stronger general agent framework, but a diagnostic decomposition that makes repair and refusal behavior observable in a constrained advisory setting.

\subsection{Safety-Sensitive Advisory Settings}

Safety-sensitive advisory settings change the cost of unsupported generation. In ordinary open-domain QA, an unsupported answer is primarily a factuality problem. In advisory contexts, unsupported or overconfident advice can also become a risk-management problem, especially when the user asks for escalation, continuation under symptoms, or precise guarantees that the evidence does not support. Recent work on retrieval-augmented language models studies whether such systems know when not to answer and how over-refusal can be steered \cite{zhou2025ralmknow,maskey2025overrefusal}; retrieval-enabled LLMs can also introduce new safety concerns rather than automatically reducing them \cite{yu2025safetydegradation}. Medical multi-agent work such as MDAgents explores adaptive collaboration for medical decision-making \cite{kim2024mdagents}, but our work is deliberately narrower: it is not a clinical system, does not evaluate patient outcomes, and does not claim deployment-ready safety.

This paper therefore treats safety-sensitive advisory RAG as a workflow-diagnosis problem. The studied system must show not only when it can answer, but also when it should produce a bounded answer, repair citation or grounding defects, or refuse. The resulting research gap is specific: prior work provides retrieval, citation, critique, repair, and agentic decomposition, but less often evaluates refusal, repair, evidence insufficiency, and safety-boundary behavior as explicit states in the same traceable advisory workflow.
